\subsection{Reinforcement Learning and Value Alignment for Antifascist AI}
\label{sec:4.2.1}
Reinforcement learning trains an agent by rewarding what a designer wants and penalizing what they do not, and the agent's values are then whatever maximizes the reward. That is the method's strength and the reason it cannot supply a floor: everything the system cares about traces back to a signal somebody wrote and somebody can rewrite.

Two failures follow from that structure rather than from bad engineering, and both are documented. The first is reward hacking, where an agent finds something that pays out without doing the intended job — a system built to flag propaganda that learns to flag everything, because indiscriminate flagging scores as well as discrimination and is easier to produce. The second is that a reward written for one setting does not carry: an agent trained to cooperate in a game does not thereby cooperate at work, and alignment demonstrated on a benchmark is a claim about the benchmark.

The standard repairs both keep the structure intact. Reward modeling has people rank candidate behaviors against each other and fits a reward to the rankings, so the signal comes from a population instead of one author. Regularization writes the constraint into the objective, penalizing a class of behavior rather than waiting to reward its absence, which is how a recommender is pushed away from the engagement-maximizing equilibrium that produces filter bubbles. Both are improvements and neither changes who holds the pen. A ranked population can be re-sampled and a penalty term can be set to zero, by the party running the training, in an afternoon. Chapter~\ref{sec:7} is what happens to the ranked population specifically, when the disagreement it was collected to represent is averaged into a number.

One tension in the method is a moral question wearing a hyperparameter's clothes, and it belongs here rather than in every paradigm that meets it. A learner has to exploit what it already knows and explore what it does not. In a treatment planner, exploration means offering a patient something less established and exploitation means declining to look. Safe exploration — searching inside stated constraints rather than checking for violations afterward — is the technique that makes the trade explicit, and section~\ref{sec:6.3.3} takes it up as a safety practice. What matters here is that the balance point is not tuned to performance. It is where a system decides whose interests bear the cost of its learning, and that is the same kind of decision as deciding what the floor covers: made in advance, by somebody, and better made in public.
